\textbf{Phần 1. Trắc nghiệm lựa chọn}
\setcounterpageref{ex}{0}

\begin{ex} %[3H13N67]
Đồng hồ điện tử hiển thị $16:45$. Lúc đó là mấy giờ chiều?
\choice
{$6$ giờ $45$ phút chiều}
{$16$ giờ $4$ phút chiều}
{\True $4$ giờ $45$ phút chiều (hay $5$ giờ kém $15$)}
{$5$ giờ $45$ phút chiều}
\loigiai{
$16:45$ là $4$ giờ $45$ phút chiều (hay 5 giờ kém 15).\\
\textbf{Đáp án của câu hỏi là C}
}
\end{ex}

\begin{ex} %[3H13N67]
Đồng hồ điện tử hiển thị $20:30$. Lúc đó là
\choice
{$10$ giờ $30$ phút tối}
{\True $8$ giờ $30$ phút tối (hay $8$ giờ rưỡi tối)}
{$8$ giờ $20$ phút tối}
{$9$ giờ rưỡi tối}
\loigiai{
$20:30$ là $8$ giờ $30$ phút tối.\\
\textbf{Đáp án của câu hỏi là B}
}
\end{ex}

\begin{ex} %[3H13H67]
Mỗi khoảng cách giữa hai vạch nhỏ liên tiếp trên mặt đồng hồ tương ứng với
\choice
{\True $1$ phút}
{$5$ phút}
{$10$ phút}
{$1$ giây}
\loigiai{
Mỗi vạch nhỏ là 1 phút.\\
\textbf{Đáp án của câu hỏi là A}
}
\end{ex}

\begin{ex} %[3H13H67]
Đồng hồ chỉ $7$ giờ $18$ phút: kim dài đang chỉ ở
\choice
{Đúng số 3}
{\True Sau số 3 thêm 3 vạch nhỏ}
{Trước số 4 hai vạch}
{Đúng số 4}
\loigiai{
Số 3 là 15 phút, thêm 3 vạch nhỏ là 18 phút.\\
\textbf{Đáp án của câu hỏi là B}
}
\end{ex}

\begin{ex} %[3H13H67]
Nếu thứ Bảy tuần này là ngày 10 thì thứ Bảy tuần sau là ngày
\choice
{$16$}
{$18$}
{$19$}
{\True $17$}
\loigiai{
$10 + 7 = 17$.\\
\textbf{Đáp án của câu hỏi là D}
}
\end{ex}

\begin{ex} %[3H13H67]
Nếu Chủ nhật tuần này là ngày 15 thì Chủ nhật tuần trước là ngày
\choice
{\True $8$}
{$7$}
{$9$}
{$10$}
\loigiai{
$15 - 7 = 8$.\\
\textbf{Đáp án của câu hỏi là A}
}
\end{ex}

\begin{ex} %[3H13V67]
Chuyến tàu khởi hành lúc 6 giờ 15 phút và đến nơi lúc 8 giờ 45 phút cùng ngày. Thời gian đi của chuyến tàu là
\choice
{$2$ giờ $15$ phút}
{$3$ giờ}
{$2$ giờ $45$ phút}
{\True $2$ giờ $30$ phút (hay $2$ giờ rưỡi)}
\loigiai{
$8$ giờ $45$ phút $- 6$ giờ $15$ phút $= 2$ giờ $30$ phút.\\
\textbf{Đáp án của câu hỏi là D}
}
\end{ex}

\begin{ex} %[3H13V67]
Bạn Nam xem một bộ phim hoạt hình từ 19 giờ 30 phút đến 20 giờ 15 phút. Bộ phim kéo dài bao nhiêu phút?
\choice
{$40$ phút}
{$50$ phút}
{\True $45$ phút}
{$35$ phút}
\loigiai{
Từ 19:30 đến 20:15 là 45 phút.\\
\textbf{Đáp án của câu hỏi là C}
}
\end{ex}

\begin{ex} %[3H13V67]
Tháng 10 có ngày 20 tháng 10 là ngày Phụ nữ Việt Nam. Nếu ngày 20 tháng 10 là Chủ nhật thì ngày 27 tháng 10 là
\choice
{$Thứ Bảy$}
{\True $Chủ nhật$}
{$Thứ Hai$}
{$Thứ Sáu$}
\loigiai{
$20 + 7 = 27$ nên vẫn là Chủ nhật.\\
\textbf{Đáp án của câu hỏi là B}
}
\end{ex}

\begin{ex} %[3H13V67]
Ngày Quốc tế Thiếu nhi 1 tháng 6 rơi vào thứ Sáu. Hỏi ngày 15 tháng 6 cùng năm đó là thứ mấy?
\choice
{\True $Thứ Sáu$}
{$Thứ Năm$}
{$Thứ Bảy$}
{$Chủ nhật$}
\loigiai{
$1 + 14 = 15$ (sau đúng 2 tuần) nên vẫn là thứ Sáu.\\
\textbf{Đáp án của câu hỏi là A}
}
\end{ex}

\begin{ex} %[3H13V67]
Bác thợ làm việc từ 7 giờ 30 phút sáng đến 11 giờ 30 phút trưa. Bác thợ đã làm việc trong bao nhiêu giờ?
\choice
{$3$ giờ}
{$5$ giờ}
{\True $4$ giờ}
{$6$ giờ}
\loigiai{
$11$ giờ $30$ phút $- 7$ giờ $30$ phút $= 4$ giờ.\\
\textbf{Đáp án của câu hỏi là C}
}
\end{ex}

\begin{ex} %[3H13C67]
Một tháng nào đó có 3 ngày thứ Năm đều là ngày chẵn. Hỏi ngày 26 của tháng đó là thứ mấy?
\choice
{$Thứ Tư$}
{$Chủ nhật$}
{$Thứ Sáu$}
{\True $Thứ Bảy$}
\loigiai{
Tháng có 3 ngày thứ Năm là ngày chẵn thì các ngày thứ Năm đó là: ngày 2, ngày 16, ngày 30.\\
Nếu ngày 30 là Thứ Năm thì ngày 26 là Chủ nhật tuần trước.\\
Ngày 30 (Thứ Năm), ngày 29 (Thứ Tư), ngày 28 (Thứ Ba), ngày 27 (Thứ Hai), ngày 26 (Chủ nhật).\\
Hoặc ngày 16 (Thứ Năm) $\rightarrow$ ngày 23 (Thứ Năm) $\rightarrow$ ngày 26 (Chủ nhật / Thứ Bảy).\\
Nếu ngày 23 là thứ Năm $\rightarrow$ 24 (T6), 25 (T7), 26 (Chủ nhật).\\
\textbf{Đáp án của câu hỏi là D}
}
\end{ex}

\textbf{Phần 2. Trắc nghiệm đúng sai}
\setcounterpageref{ex}{0}

\begin{ex} %[3H13H67]
Xét tính đúng sai về cách đọc giờ và lịch:
\choiceTF
{\True $14$ giờ $30$ phút tức là $2$ giờ $30$ phút chiều.}
{\True Trong một tháng, hai ngày cùng thứ liền nhau cách nhau 7 ngày.}
{Tháng 2 luôn có 30 ngày.}
{\True Tháng 2 chỉ có 28 hoặc 29 ngày.}
\loigiai{
\begin{itemize}
    \item $14:30$ là 2 giờ rưỡi chiều: Đúng.
    \item Hai ngày cùng thứ cách nhau 7 ngày: Đúng.
    \item Tháng 2 không bao giờ có 30 ngày. Ý 3 Sai, Ý 4 Đúng.
\end{itemize}
\begin{itemize}
    \item a) \textbf{Đúng}.
    \item b) \textbf{Đúng}.
    \item c) \textbf{Sai}.
    \item d) \textbf{Đúng}.
\end{itemize}
}
\end{ex}

\begin{ex} %[3H13V67]
Một hội trại diễn ra từ 8 giờ sáng thứ Bảy đến 16 giờ chiều Chủ nhật tuần đó. Xét tính đúng sai:
\choiceTF
{\True Từ 8 giờ sáng thứ Bảy đến 8 giờ sáng Chủ nhật là 24 giờ.}
{\True Từ 8 giờ sáng Chủ nhật đến 16 giờ chiều Chủ nhật là 8 giờ.}
{Tổng thời gian diễn ra hội trại là 30 giờ.}
{\True Tổng thời gian diễn ra hội trại là 32 giờ ($24 + 8 = 32$).}
\loigiai{
\begin{itemize}
    \item 1 ngày đêm là 24 giờ: Đúng.
    \item $16 - 8 = 8$ giờ: Đúng.
    \item Tổng: $24 + 8 = 32$ giờ $\neq 30$ giờ. Ý 3 Sai, Ý 4 Đúng.
\end{itemize}
\begin{itemize}
    \item a) \textbf{Đúng}.
    \item b) \textbf{Đúng}.
    \item c) \textbf{Sai}.
    \item d) \textbf{Đúng}.
\end{itemize}
}
\end{ex}

\textbf{Phần 3. Trả lời ngắn}
\setcounterpageref{ex}{0}

\begin{ex} %[3H13H67]
Đồng hồ $17:00$ tương ứng với mấy giờ chiều?
\shortans{5}
\loigiai{
\textbf{Đáp số:} $5$\\
5 giờ chiều.
}
\end{ex}

\begin{ex} %[3H13H67]
Nếu thứ Ba là ngày 5 thì thứ Ba tuần sau là ngày mấy?
\shortans{12}
\loigiai{
\textbf{Đáp số:} $12$\\
$5 + 7 = 12$.
}
\end{ex}

\begin{ex} %[3H13V67]
Một trận bóng đá bắt đầu lúc 15 giờ 30 phút và kết thúc lúc 17 giờ 15 phút. Trận đấu diễn ra trong bao nhiêu phút?
\shortans{105}
\loigiai{
\textbf{Đáp số:} $105$\\
1 giờ 45 phút $= 105$ phút.
}
\end{ex}

\begin{ex} %[3H13V67]
Nếu ngày 25 tháng 12 là thứ Tư thì ngày 18 tháng 12 cùng năm là thứ mấy (ghi số thứ, ví dụ thứ Tư ghi 4)?
\shortans{4}
\loigiai{
\textbf{Đáp số:} $4$\\
$25 - 7 = 18$ nên vẫn là thứ Tư (ghi 4).
}
\end{ex}

\begin{ex} %[3H13V67]
Từ 14 giờ đến 17 giờ 30 phút có tất cả bao nhiêu phút?
\shortans{210}
\loigiai{
\textbf{Đáp số:} $210$\\
3 giờ 30 phút $= 180 + 30 = 210$ phút.
}
\end{ex}

\begin{ex} %[3H13C67]
Tháng 3 có 31 ngày. Biết ngày 1 tháng 3 là thứ Sáu. Hỏi tháng 3 có bao nhiêu ngày Chủ nhật?
\shortans{5}
\loigiai{
\textbf{Đáp số:} $5$\\
Các ngày Chủ nhật là ngày: 3, 10, 17, 24, 31 (có 5 ngày Chủ nhật).
}
\end{ex}
